\chapter{Física atómica}
\label{cap:fisica_atomica}

Este capítulo aborda la estructura y el comportamiento de los átomos desde las perspectivas histórica y cuántica. Se estudian los espectros atómicos de gases, el modelo atómico clásico e inestable, la cuantización en el modelo del átomo de hidrógeno de Bohr, el tratamiento mecánico-cuántico tridimensional con las correspondientes funciones de onda e interpretación cuántica de números cuánticos (incluido el espín), el principio de exclusión de Pauli que rige el llenado de la tabla periódica, los espectros de rayos X de capas internas (Moseley) y los fundamentos de la emisión estimulada en láseres.

\section*{Piense, dialogue y comparta}
\begin{enumerate}
    \subimport{piense_01/}{ejercicio.tex}
    \subimport{piense_02/}{ejercicio.tex}
\end{enumerate}

\section*{Problemas}

\subsection*{SECCIÓN 41.1 Espectros atómicos de los gases}
\begin{enumerate}
    \subimport{prob_41_1_01/}{ejercicio.tex}
    \subimport{prob_41_1_02/}{ejercicio.tex}
    \subimport{prob_41_1_03/}{ejercicio.tex}
\end{enumerate}

\subsection*{SECCIÓN 41.2 Los primeros modelos del átomo}
\begin{enumerate}
    \setcounter{enumi}{3}
    \subimport{prob_41_2_04/}{ejercicio.tex}
\end{enumerate}

\subsection*{SECCIÓN 41.3 Modelo del átomo de Bohr de hidrógeno}
\begin{enumerate}
    \setcounter{enumi}{4}
    \subimport{prob_41_3_05/}{ejercicio.tex}
    \subimport{prob_41_3_06/}{ejercicio.tex}
    \subimport{prob_41_3_07/}{ejercicio.tex}
    \subimport{prob_41_3_08/}{ejercicio.tex}
    \subimport{prob_41_3_09/}{ejercicio.tex}
    \subimport{prob_41_3_10/}{ejercicio.tex}
    \subimport{prob_41_3_11/}{ejercicio.tex}
    \subimport{prob_41_3_12/}{ejercicio.tex}
    \subimport{prob_41_3_13/}{ejercicio.tex}
\end{enumerate}

\subsection*{SECCIÓN 41.4 Modelo cuántico del átomo de hidrógeno}
\begin{enumerate}
    \setcounter{enumi}{13}
    \subimport{prob_41_4_14/}{ejercicio.tex}
\end{enumerate}

\subsection*{SECCIÓN 41.5 Las funciones de onda para el hidrógeno}
\begin{enumerate}
    \setcounter{enumi}{14}
    \subimport{prob_41_5_15/}{ejercicio.tex}
    \subimport{prob_41_5_16/}{ejercicio.tex}
    \subimport{prob_41_5_17/}{ejercicio.tex}
\end{enumerate}

\subsection*{SECCIÓN 41.6 Interpretación física de los números cuánticos}
\begin{enumerate}
    \setcounter{enumi}{17}
    \subimport{prob_41_6_18/}{ejercicio.tex}
    \subimport{prob_41_6_19/}{ejercicio.tex}
    \subimport{prob_41_6_20/}{ejercicio.tex}
    \subimport{prob_41_6_21/}{ejercicio.tex}
    \subimport{prob_41_6_22/}{ejercicio.tex}
    \subimport{prob_41_6_23/}{ejercicio.tex}
    \subimport{prob_41_6_24/}{ejercicio.tex}
\end{enumerate}

\subsection*{SECCIÓN 41.7 El principio de exclusión y la tabla periódica}
\begin{enumerate}
    \setcounter{enumi}{24}
    \subimport{prob_41_7_25/}{ejercicio.tex}
    \subimport{prob_41_7_26/}{ejercicio.tex}
    \subimport{prob_41_7_27/}{ejercicio.tex}
    \subimport{prob_41_7_28/}{ejercicio.tex}
    \subimport{prob_41_7_29/}{ejercicio.tex}
\end{enumerate}

\subsection*{SECCIÓN 41.8 Más sobre los espectros atómicos: el visible y el rayo X}
\begin{enumerate}
    \setcounter{enumi}{29}
    \subimport{prob_41_8_30/}{ejercicio.tex}
    \subimport{prob_41_8_31/}{ejercicio.tex}
    \subimport{prob_41_8_32/}{ejercicio.tex}
    \subimport{prob_41_8_33/}{ejercicio.tex}
    \subimport{prob_41_8_34/}{ejercicio.tex}
\end{enumerate}

\subsection*{SECCIÓN 41.10 Láseres}
\begin{enumerate}
    \setcounter{enumi}{34}
    \subimport{prob_41_10_35/}{ejercicio.tex}
    \subimport{prob_41_10_36/}{ejercicio.tex}
\end{enumerate}

\section*{PROBLEMAS ADICIONALES}
\begin{enumerate}
    \setcounter{enumi}{36}
    \subimport{prob_adicional_37/}{ejercicio.tex}
    \subimport{prob_adicional_38/}{ejercicio.tex}
    \subimport{prob_adicional_39/}{ejercicio.tex}
    \subimport{prob_adicional_40/}{ejercicio.tex}
    \subimport{prob_adicional_41/}{ejercicio.tex}
    \subimport{prob_adicional_42/}{ejercicio.tex}
    \subimport{prob_adicional_43/}{ejercicio.tex}
    \subimport{prob_adicional_44/}{ejercicio.tex}
    \subimport{prob_adicional_45/}{ejercicio.tex}
    \subimport{prob_adicional_46/}{ejercicio.tex}
    \subimport{prob_adicional_47/}{ejercicio.tex}
    \subimport{prob_adicional_48/}{ejercicio.tex}
    \subimport{prob_adicional_49/}{ejercicio.tex}
    \subimport{prob_adicional_50/}{ejercicio.tex}
\end{enumerate}

\section*{PROBLEMAS DE DESAFÍO}
\begin{enumerate}
    \setcounter{enumi}{50}
    \subimport{prob_desafio_51/}{ejercicio.tex}
    \subimport{prob_desafio_52/}{ejercicio.tex}
\end{enumerate}
